\section{The Nine Faces of Pseudo-Needs}

\subsection{Definition}
A \textbf{strategy} is the specific way, method, or demand we adopt to
satisfy a certain need.
It is specific, personal, and sometimes time-sensitive \citep{gordon1970}.
A \textbf{need} is the deeper, universal internal desire behind the strategy.

For example:
\begin{itemize}
    \item \textbf{Strategy:} ``You must be home by 10 PM every night.''
    \item \textbf{Need:} Security, companionship, connection \citep{kabatzinn1990}.
\end{itemize}

\begin{itemize}
    \item \textbf{Strategy:} ``You have to message me every day saying you
    love me.''
    \item \textbf{Need:} To be valued, to be loved, confirmation \citep{brown2012}.
\end{itemize}

\begin{itemize}
    \item \textbf{Strategy:} ``You cannot eat alone with that colleague.''
    \item \textbf{Need:} Security, trust, loyalty \citep{kohn2005}.
\end{itemize}

When we mistake a strategy for a need, we begin to use various means
(demands, pressure, judgment, punishment) to force the other party to
execute the ``only correct way of satisfaction.''
This often triggers the other party's defense, creating new conflicts,
while the real need remains unmet \citep{faber1980}.

\paragraph{Case Analysis: Confusing Needs and Strategies}

The wife asks her husband to ``come home early every day.''
The husband considers this unreasonable control and feels oppressed and
unfree.
Both sides argue, and the conflict escalates.

But if we peel off the shell of the strategy and explore the needs behind
it:
\begin{itemize}
    \item \textbf{Wife's Needs:} Companionship, connection, security,
    feeling valued in the family \citep{cohen2001}.
    \item \textbf{Husband's Needs:} Autonomy, respect, feeling valued at
    work.
\end{itemize}

These two needs do not conflict!
What conflicts are the strategies---``coming home early'' vs. ``completing
work tasks.''
When both sides see each other's needs clearly, they can jointly find
creative strategies that meet both people's needs simultaneously:
For example, communicating the reasons and times for returning late in
advance, or agreeing on exclusive companionship time \citep{greene1998}.
\textbf{Seeing needs, rather than obsessing over strategies, is the key to
transforming conflict into connection.}

\subsection{Pseudo-Needs}
\begin{table}[h]
\centering
\caption{Common Pseudo-Needs vs. Real Needs}
\label{tab:pseudo_needs}
\begin{tabular}{|p{4.5cm}|p{3.5cm}|p{6.5cm}|}
\hline
\textbf{Pseudo-Need Expression (Strategy/Demand)} & \textbf{Underlying Real Need} & \textbf{Possible Communication Risk} \\ \hline
``I need you to want me every day.'' & Connection, being valued, security & The other person feels forced to execute a task, rather than being emotionally desired. \\ \hline
``I need you to be more responsible.'' & Support, ease, shared burden & Implies personality judgment, highly likely to trigger defense and counterattack. \\ \hline
``I need you to listen to me.'' & Being seen, respect, sense of efficacy & Evolves into a power struggle, ignoring the basis of equal dialogue. \\ \hline
``I need you to give me some space.'' & Autonomy, tranquility, self-regulation & Without background explanation, the other person will feel neglected or abandoned. \\ \hline
``I need you to agree with my views.'' & Being understood, confirmation, harmony & Forcing unified concepts, leading the other person to suppress real thoughts. \\ \hline
``I need you to change your attitude.'' & Respect, security, acceptance & Expression is too vague, the other person feels frustrated due to not knowing what to do. \\ \hline
``I need you to stop being so selfish.'' & Being considered, fairness, mutuality & Pure moral accusation, completely blocking the exploration of needs. \\ \hline
``I need you to share the load.'' & Support, collaboration, fairness & If issued as an order, it weakens the motivation for voluntary contribution. \\ \hline
``I need you to trust me.'' & Connection, acceptance, honesty & Trust is the result of interaction and cannot be obtained through a unilateral demand. \\ \hline
\end{tabular}
\end{table}

As shown in Table \ref{tab:pseudo_needs}, similar to ``pseudo-feelings,''
we also have expressions of ``pseudo-needs''---those words that seem to be
stating needs but are actually expressing judgments or strategic demands.
Identifying pseudo-needs and translating them into expressions of real
needs can fundamentally change the texture of communication---it moves the
focus from ``your behavior'' to ``my heart,'' creating more possibilities
for both sides \citep{siegel2011}.

\paragraph{Psychological Mechanism: Why Do We Habitually Express Pseudo-Needs?}
The reason we often fall into the expression pattern of pseudo-needs usually
originates from deep psychological defense mechanisms and the evasion of
responsibility:
\begin{itemize}
    \item \textbf{Defense Mechanisms and False Sense of Control:} Expressing
    real needs (such as ``I feel insecure and need your support'') makes us
    appear vulnerable, as if handing over the ``power to hurt ourselves''
    to the other party \citep{lazarus1991}.
    In contrast, expressing pseudo-needs (such as ``I need you to be more
    responsible'') is more like issuing an order.
    This condescending posture can bring a false ``sense of control,''
    allowing us to hide behind tough language and avoid facing our inner
    powerlessness.
    \item \textbf{Transferring Responsibility:} Pseudo-needs are essentially
    demanding that the other party be responsible for our feelings.
    This allows us to escape the responsibility of exploring our own needs,
    and instead push the burden of solving the problem entirely onto the
    other party.
    However, this transfer usually triggers the other party's defense
    (such as counterattack or silence), causing the problem to fall into a
    stalemate.
    \item \textbf{The Influence of Culture and Socialization:} From childhood
    to adulthood, we may be educated to be ``good children'' or ``good
    partners,'' instilled with concepts of ``you should...'' and ``you
    must...''.
    This cultural background reinforces our habit of using judgments and
    demands to express needs, rather than directly expressing inner desires.
    \item \textbf{Default Status Differences:} In relationships such as
    parent and child, superior and subordinate, we may be accustomed to
    using imperative language to express needs because it conforms to our
    default status difference in the relationship.
\end{itemize}

When we cannot clearly identify the universal needs deep within our hearts,
these needs often ``deform'' into aggressive or controlling pseudo-needs
(strategies).

\begin{keyinsight}
Confusing our specific strategies for getting what we want with the actual
universal needs that drive us is a primary source of conflict.
When we mistakenly frame demands or behavioral controls as our ``needs,''
we inevitably trigger defensiveness; instead, we must learn to articulate
our true underlying vulnerabilities and desires.
\end{keyinsight}
